% Reconstructed text from V5 PDF page 10; non-native source
\par\smallskip\noindent\rule{\textwidth}{0.4pt}\par{\sffamily\footnotesize Original V5 page 10 -- recovered text}\par\begin{flushleft}\footnotesize\ttfamily
\par\smallskip
\par\smallskip
\noindent Diagnostics and Serving Infrastructure for MAS 10\par
\par\smallskip
\par\smallskip
\par\smallskip
\noindent finding to its source section, table, or figure; and (10) cross-layer analysis notes assessing the\par
\noindent system’s relevance to the diagnostic-infrastructure interaction.\par
\noindent  This structured template ensures that the comparison tables in Sections 4 and 5 report con-\par
\noindent sistent dimensions across all systems. The template was developed iteratively: the first five\par
\noindent systems (SGLang, ScaleSim, DistServe, Mooncake, Parrot) were extracted using a preliminary\par
\noindent template, and the template was refined based on which dimensions proved most informative for\par
\noindent cross-system comparison. All subsequent extractions used the refined template.\par
\noindent  A verification step was applied to each extraction: quantitative claims in the extraction tem-\par
\noindent plate were cross-checked against the source paper’s tables and figures (not just the text). Several\par
\noindent discrepancies were identified and corrected: for example, a preliminary extraction attributed\par
\noindent 66\% accuracy to FAMAS, but the source paper’s Table 2 shows 62.07\% as the best agent-level\par
\noindent figure on handcrafted scenarios; similarly, TrajAudit’s trajectory folding was initially recorded\par
\noindent as 86–92\% but the source paper (§VII-D) reports 94.6\%.\par
\par\smallskip
\noindent 3.6 Corpus Summary\par
\noindent The 60 core table entries comprise 56 entries assessed at Tier 1 (T1a or T1b) and four at\par
\noindent Tier 2. The core failure-localization table includes eight full-text-verified (T1) methods and\par
\noindent four abstract-verified (T2) methods. The T2 entries — StepFinder, FALAT, DoVer, and Veri-\par
\noindent fyMAS — are retained for landscape coverage but are not used to support fine-grained mech-\par
\noindent anistic or quantitative cross-method comparisons unless their primary full texts are indepen-\par
\noindent dently checked. Additional contextual works discussed in the running text are not included in\par
\noindent this count. By category: 12 failure-localization methods (8 T1, 4 T2), 6 observability entries\par
\noindent (all T1), 8 safety and guardrail systems (4 T1a, 4 T1b), 22 infrastructure systems (all T1), and\par
\noindent 12 evaluation benchmarks (all T1). See Appendix A for the complete census. Additionally, we\par
\noindent include a retention-ablation experiment on real TraceElephant traces (220 failures) to support\par
\noindent the validation protocol proposed in Section 7.\par
\noindent  Figure 2 illustrates the corpus assembly process.\par
\par\smallskip
\par\smallskip
\noindent 4 Diagnostics\par
\noindent This section organizes the diagnostic landscape along functional roles rather than implementa-\par
\noindent tion categories, following the taxonomy introduced in Section 3.4. We cover failure localization\par
\noindent (§4.1), observability and tracing (§4.2), safety intervention and guardrails (§4.3), and diagnos-\par
\noindent tic benchmarks (§4.4). For each subsection, methods are characterized along the six method\par
\noindent dimensions defined in Section 3.4: evidence inputs, training requirements, replay demands, at-\par
\noindent tribution granularity, intervention timing, and runtime cost. Figure 3 provides a matrix view of\par
\noindent the taxonomy with representative systems placed in cells.\par
\par\smallskip
\noindent 4.1 Failure Localization\par
\noindent Failure localization identifies which agent, step, or action caused an observed failure. We iden-\par
\noindent tify four methodological families among the 8 Tier-1 methods reviewed: statistical methods\par
\noindent adapted from software fault localization, LLM-guided methods that use a language model as\par
\noindent an investigator, RL-trained methods that learn attribution policies, and training-free signal-\par
\noindent extraction methods that exploit model internals.\par
\end{flushleft}
